\IfFileExists{stacks-project.cls}{%
\documentclass{stacks-project}
}{%
\documentclass{amsart}
}

% For dealing with references we use the comment environment
\usepackage{verbatim}
\newenvironment{reference}{\comment}{\endcomment}
%\newenvironment{reference}{}{}
\newenvironment{slogan}{\comment}{\endcomment}
\newenvironment{history}{\comment}{\endcomment}

% For commutative diagrams we use Xy-pic
\usepackage[all]{xy}

% We use 2cell for 2-commutative diagrams.
\xyoption{2cell}
\UseAllTwocells

% We use multicol for the list of chapters between chapters
\usepackage{multicol}

% This is generally recommended for better output
\usepackage{lmodern}
\usepackage[T1]{fontenc}

% For cross-file-references
\usepackage{xr-hyper}

% Package for hypertext links:
\usepackage{hyperref}

% For any local file, say "hello.tex" you want to link to please
% use \externaldocument[hello-]{hello}
\externaldocument[introduction-]{introduction}
\externaldocument[conventions-]{conventions}
\externaldocument[sets-]{sets}
\externaldocument[categories-]{categories}
\externaldocument[topology-]{topology}
\externaldocument[sheaves-]{sheaves}
\externaldocument[sites-]{sites}
\externaldocument[stacks-]{stacks}
\externaldocument[fields-]{fields}
\externaldocument[algebra-]{algebra}
\externaldocument[brauer-]{brauer}
\externaldocument[homology-]{homology}
\externaldocument[derived-]{derived}
\externaldocument[simplicial-]{simplicial}
\externaldocument[more-algebra-]{more-algebra}
\externaldocument[smoothing-]{smoothing}
\externaldocument[modules-]{modules}
\externaldocument[sites-modules-]{sites-modules}
\externaldocument[injectives-]{injectives}
\externaldocument[cohomology-]{cohomology}
\externaldocument[sites-cohomology-]{sites-cohomology}
\externaldocument[dga-]{dga}
\externaldocument[dpa-]{dpa}
\externaldocument[sdga-]{sdga}
\externaldocument[hypercovering-]{hypercovering}
\externaldocument[schemes-]{schemes}
\externaldocument[constructions-]{constructions}
\externaldocument[properties-]{properties}
\externaldocument[morphisms-]{morphisms}
\externaldocument[coherent-]{coherent}
\externaldocument[divisors-]{divisors}
\externaldocument[limits-]{limits}
\externaldocument[varieties-]{varieties}
\externaldocument[topologies-]{topologies}
\externaldocument[descent-]{descent}
\externaldocument[perfect-]{perfect}
\externaldocument[more-morphisms-]{more-morphisms}
\externaldocument[flat-]{flat}
\externaldocument[groupoids-]{groupoids}
\externaldocument[more-groupoids-]{more-groupoids}
\externaldocument[etale-]{etale}
\externaldocument[chow-]{chow}
\externaldocument[intersection-]{intersection}
\externaldocument[pic-]{pic}
\externaldocument[weil-]{weil}
\externaldocument[adequate-]{adequate}
\externaldocument[dualizing-]{dualizing}
\externaldocument[duality-]{duality}
\externaldocument[discriminant-]{discriminant}
\externaldocument[derham-]{derham}
\externaldocument[local-cohomology-]{local-cohomology}
\externaldocument[algebraization-]{algebraization}
\externaldocument[curves-]{curves}
\externaldocument[resolve-]{resolve}
\externaldocument[models-]{models}
\externaldocument[functors-]{functors}
\externaldocument[equiv-]{equiv}
\externaldocument[pione-]{pione}
\externaldocument[etale-cohomology-]{etale-cohomology}
\externaldocument[proetale-]{proetale}
\externaldocument[relative-cycles-]{relative-cycles}
\externaldocument[more-etale-]{more-etale}
\externaldocument[trace-]{trace}
\externaldocument[crystalline-]{crystalline}
\externaldocument[spaces-]{spaces}
\externaldocument[spaces-properties-]{spaces-properties}
\externaldocument[spaces-morphisms-]{spaces-morphisms}
\externaldocument[decent-spaces-]{decent-spaces}
\externaldocument[spaces-cohomology-]{spaces-cohomology}
\externaldocument[spaces-limits-]{spaces-limits}
\externaldocument[spaces-divisors-]{spaces-divisors}
\externaldocument[spaces-over-fields-]{spaces-over-fields}
\externaldocument[spaces-topologies-]{spaces-topologies}
\externaldocument[spaces-descent-]{spaces-descent}
\externaldocument[spaces-perfect-]{spaces-perfect}
\externaldocument[spaces-more-morphisms-]{spaces-more-morphisms}
\externaldocument[spaces-flat-]{spaces-flat}
\externaldocument[spaces-groupoids-]{spaces-groupoids}
\externaldocument[spaces-more-groupoids-]{spaces-more-groupoids}
\externaldocument[bootstrap-]{bootstrap}
\externaldocument[spaces-pushouts-]{spaces-pushouts}
\externaldocument[spaces-chow-]{spaces-chow}
\externaldocument[groupoids-quotients-]{groupoids-quotients}
\externaldocument[spaces-more-cohomology-]{spaces-more-cohomology}
\externaldocument[spaces-simplicial-]{spaces-simplicial}
\externaldocument[spaces-duality-]{spaces-duality}
\externaldocument[formal-spaces-]{formal-spaces}
\externaldocument[restricted-]{restricted}
\externaldocument[spaces-resolve-]{spaces-resolve}
\externaldocument[formal-defos-]{formal-defos}
\externaldocument[defos-]{defos}
\externaldocument[cotangent-]{cotangent}
\externaldocument[examples-defos-]{examples-defos}
\externaldocument[algebraic-]{algebraic}
\externaldocument[examples-stacks-]{examples-stacks}
\externaldocument[stacks-sheaves-]{stacks-sheaves}
\externaldocument[criteria-]{criteria}
\externaldocument[artin-]{artin}
\externaldocument[quot-]{quot}
\externaldocument[stacks-properties-]{stacks-properties}
\externaldocument[stacks-morphisms-]{stacks-morphisms}
\externaldocument[stacks-limits-]{stacks-limits}
\externaldocument[stacks-cohomology-]{stacks-cohomology}
\externaldocument[stacks-perfect-]{stacks-perfect}
\externaldocument[stacks-introduction-]{stacks-introduction}
\externaldocument[stacks-more-morphisms-]{stacks-more-morphisms}
\externaldocument[stacks-geometry-]{stacks-geometry}
\externaldocument[moduli-]{moduli}
\externaldocument[moduli-curves-]{moduli-curves}
\externaldocument[examples-]{examples}
\externaldocument[exercises-]{exercises}
\externaldocument[guide-]{guide}
\externaldocument[desirables-]{desirables}
\externaldocument[coding-]{coding}
\externaldocument[obsolete-]{obsolete}
\externaldocument[fdl-]{fdl}
\externaldocument[index-]{index}

% Theorem environments.
%
\theoremstyle{plain}
\newtheorem{theorem}[subsection]{Theorem}
\newtheorem{proposition}[subsection]{Proposition}
\newtheorem{lemma}[subsection]{Lemma}

\theoremstyle{definition}
\newtheorem{definition}[subsection]{Definition}
\newtheorem{example}[subsection]{Example}
\newtheorem{exercise}[subsection]{Exercise}
\newtheorem{situation}[subsection]{Situation}

\theoremstyle{remark}
\newtheorem{remark}[subsection]{Remark}
\newtheorem{remarks}[subsection]{Remarks}

\numberwithin{equation}{subsection}

% Macros
%
\def\lim{\mathop{\mathrm{lim}}\nolimits}
\def\colim{\mathop{\mathrm{colim}}\nolimits}
\def\Spec{\mathop{\mathrm{Spec}}}
\def\Hom{\mathop{\mathrm{Hom}}\nolimits}
\def\Ext{\mathop{\mathrm{Ext}}\nolimits}
\def\SheafHom{\mathop{\mathcal{H}\!\mathit{om}}\nolimits}
\def\SheafExt{\mathop{\mathcal{E}\!\mathit{xt}}\nolimits}
\def\Sch{\mathit{Sch}}
\def\Mor{\mathop{\mathrm{Mor}}\nolimits}
\def\Ob{\mathop{\mathrm{Ob}}\nolimits}
\def\Sh{\mathop{\mathit{Sh}}\nolimits}
\def\NL{\mathop{N\!L}\nolimits}
\def\CH{\mathop{\mathrm{CH}}\nolimits}
\def\proetale{{pro\text{-}\acute{e}tale}}
\def\etale{{\acute{e}tale}}
\def\QCoh{\mathit{QCoh}}
\def\Ker{\mathop{\mathrm{Ker}}}
\def\Im{\mathop{\mathrm{Im}}}
\def\Coker{\mathop{\mathrm{Coker}}}
\def\Coim{\mathop{\mathrm{Coim}}}

% Boxtimes
%
\DeclareMathSymbol{\boxtimes}{\mathbin}{AMSa}{"02}

%
% Macros for moduli stacks/spaces
%
\def\QCohstack{\mathcal{QC}\!\mathit{oh}}
\def\Cohstack{\mathcal{C}\!\mathit{oh}}
\def\Spacesstack{\mathcal{S}\!\mathit{paces}}
\def\Quotfunctor{\mathrm{Quot}}
\def\Hilbfunctor{\mathrm{Hilb}}
\def\Curvesstack{\mathcal{C}\!\mathit{urves}}
\def\Polarizedstack{\mathcal{P}\!\mathit{olarized}}
\def\Complexesstack{\mathcal{C}\!\mathit{omplexes}}
% \Pic is the operator that assigns to X its picard group, usage \Pic(X)
% \Picardstack_{X/B} denotes the Picard stack of X over B
% \Picardfunctor_{X/B} denotes the Picard functor of X over B
\def\Pic{\mathop{\mathrm{Pic}}\nolimits}
\def\Picardstack{\mathcal{P}\!\mathit{ic}}
\def\Picardfunctor{\mathrm{Pic}}
\def\Deformationcategory{\mathcal{D}\!\mathit{ef}}
\usepackage{cleveref}
\crefname{originalreference}{original reference}{original references}
\Crefname{originalreference}{Original reference}{Original references}
\makeatletter
\newlabel{groupoids-lemma-find-invariant-affine}{{03JE}{source}{}{sourcepointer-1}{}}
\newlabel{groupoids-lemma-find-invariant-affine@cref}{{[originalreference][0][]03JE}{[1][0][]source}{}{}{}}
\makeatother
\crefname{originalreference}{original reference}{original references}
\Crefname{originalreference}{Original reference}{Original references}
\makeatletter
\newlabel{groupoids-section-invariant}{{03LN}{source}{}{sourcepointer-2}{}}
\newlabel{groupoids-section-invariant@cref}{{[originalreference][0][]03LN}{[1][0][]source}{}{}{}}
\makeatother
\crefname{originallemma}{original lemma}{original lemmas}
\Crefname{originallemma}{Original lemma}{Original lemmas}
\makeatletter
\newlabel{lemma-strong-splitting-scheme}{{04RX}{source}{}{sourcepointer-3}{}}
\newlabel{lemma-strong-splitting-scheme@cref}{{[originallemma][0][]04RX}{[1][0][]source}{}{}{}}
\makeatother
\crefname{originalreference}{original reference}{original references}
\Crefname{originalreference}{Original reference}{Original references}
\makeatletter
\newlabel{more-groupoids-lemma-restrict-preserves-type}{{04MP}{source}{}{sourcepointer-4}{}}
\newlabel{more-groupoids-lemma-restrict-preserves-type@cref}{{[originalreference][0][]04MP}{[1][0][]source}{}{}{}}
\makeatother
\crefname{originalreference}{original reference}{original references}
\Crefname{originalreference}{Original reference}{Original references}
\makeatletter
\newlabel{more-morphisms-lemma-quasi-finite-separated-quasi-affine}{{02LR}{source}{}{sourcepointer-5}{}}
\newlabel{more-morphisms-lemma-quasi-finite-separated-quasi-affine@cref}{{[originalreference][0][]02LR}{[1][0][]source}{}{}{}}
\makeatother
\crefname{originalreference}{original reference}{original references}
\Crefname{originalreference}{Original reference}{Original references}
\makeatletter
\newlabel{morphisms-lemma-etale-locally-quasi-finite}{{03WS}{source}{}{sourcepointer-6}{}}
\newlabel{morphisms-lemma-etale-locally-quasi-finite@cref}{{[originalreference][0][]03WS}{[1][0][]source}{}{}{}}
\makeatother
\crefname{originalreference}{original reference}{original references}
\Crefname{originalreference}{Original reference}{Original references}
\makeatletter
\newlabel{morphisms-lemma-finite-flat}{{02KB}{source}{}{sourcepointer-7}{}}
\newlabel{morphisms-lemma-finite-flat@cref}{{[originalreference][0][]02KB}{[1][0][]source}{}{}{}}
\makeatother
\crefname{originalreference}{original reference}{original references}
\Crefname{originalreference}{Original reference}{Original references}
\makeatletter
\newlabel{properties-lemma-ample-finite-set-in-affine}{{01ZY}{source}{}{sourcepointer-8}{}}
\newlabel{properties-lemma-ample-finite-set-in-affine@cref}{{[originalreference][0][]01ZY}{[1][0][]source}{}{}{}}
\makeatother
\crefname{originalreference}{original reference}{original references}
\Crefname{originalreference}{Original reference}{Original references}
\makeatletter
\newlabel{schemes-lemma-quasi-compact-affine}{{01K4}{source}{}{sourcepointer-9}{}}
\newlabel{schemes-lemma-quasi-compact-affine@cref}{{[originalreference][0][]01K4}{[1][0][]source}{}{}{}}
\makeatother
\crefname{originalreference}{original reference}{original references}
\Crefname{originalreference}{Original reference}{Original references}
\makeatletter
\newlabel{situation-strong-splitting}{{04RL}{source}{}{sourcepointer-10}{}}
\newlabel{situation-strong-splitting@cref}{{[originalreference][0][]04RL}{[1][0][]source}{}{}{}}
\makeatother
\title{Strong splitting on an affine base shrinks to an affine invariant open}
\author{The Stacks Project Authors}
\date{stacks-project:spaces-more-groupoids.tex:local-snapshot-20261012; GNU FDL 1.2 or later}
\begin{document}
\maketitle
\noindent\textit{Editorial note.} No omission markers in-span; the author's verification sentence for (a), (b), (c) is an equality check written out in the source.\par\medskip
\noindent\textit{Allowed standard prerequisites (editorial).}\begin{itemize}
\item Scheme-version strong splitting. lemma-strong-splitting-scheme (carded L26661). Supplies the separated etale U' -> U and the point u' with P the strong splitting.
\item Restriction preserves flatness and finite presentation. More on Groupoids, lemma restrict-preserves-type (cross-chapter, named). s', t' stay flat locally finitely presented after restriction.
\item Quasi-affineness of quasi-finite separated morphisms into affines. More on Morphisms, lemma quasi-finite-separated-quasi-affine (cross-chapter, named). Zariski's Main Theorem form making W quasi-affine.
\item Ample opens through finite sets and invariant affines. Properties, lemma ample-finite-set-in-affine; Groupoids, lemma find-invariant-affine (cross-chapter, named). Finds the affine P-invariant open U'' containing u'.
\item Finite flat is finite locally free. Morphisms, lemma finite-flat (cross-chapter, named). t(s\textasciicircum{}\{-1\}(u')) finite set of points.
\end{itemize}\medskip
\ifdefined\hypertarget\hypertarget{sourcepointer-1}{}\fi
\noindent\textit{Source reference.} Original reference 03JE: focused-seeds/geometry/S9-smg-strong-splitting-affine-scheme/ref-tags.txt, source lines 5--5. Navigation pointer only: the counted proof cites this statement by \textbackslash{}ref as a standing input; the deduction the counted proof has to make is carried out inside the counted span, and no part of that deduction is delegated to this pointer.\par\medskip
\ifdefined\hypertarget\hypertarget{sourcepointer-2}{}\fi
\noindent\textit{Source reference.} Original reference 03LN: focused-seeds/geometry/S9-smg-strong-splitting-affine-scheme/ref-tags.txt, source lines 6--6. Navigation pointer only: the counted proof cites this statement by \textbackslash{}ref as a standing input; the deduction the counted proof has to make is carried out inside the counted span, and no part of that deduction is delegated to this pointer.\par\medskip
\ifdefined\hypertarget\hypertarget{sourcepointer-3}{}\fi
\noindent\textit{Source reference.} Original lemma 04RX: sources/stacks-spaces-more-groupoids/pinned-original-source.tex, source lines 2689--2689. Navigation pointer only: the counted proof cites this statement by \textbackslash{}ref as a standing input; the deduction the counted proof has to make is carried out inside the counted span, and no part of that deduction is delegated to this pointer.\par\medskip
\ifdefined\hypertarget\hypertarget{sourcepointer-4}{}\fi
\noindent\textit{Source reference.} Original reference 04MP: focused-seeds/geometry/S9-smg-strong-splitting-affine-scheme/ref-tags.txt, source lines 7--7. Navigation pointer only: the counted proof cites this statement by \textbackslash{}ref as a standing input; the deduction the counted proof has to make is carried out inside the counted span, and no part of that deduction is delegated to this pointer.\par\medskip
\ifdefined\hypertarget\hypertarget{sourcepointer-5}{}\fi
\noindent\textit{Source reference.} Original reference 02LR: focused-seeds/geometry/S9-smg-strong-splitting-affine-scheme/ref-tags.txt, source lines 8--8. Navigation pointer only: the counted proof cites this statement by \textbackslash{}ref as a standing input; the deduction the counted proof has to make is carried out inside the counted span, and no part of that deduction is delegated to this pointer.\par\medskip
\ifdefined\hypertarget\hypertarget{sourcepointer-6}{}\fi
\noindent\textit{Source reference.} Original reference 03WS: focused-seeds/geometry/S9-smg-strong-splitting-affine-scheme/ref-tags.txt, source lines 9--9. Navigation pointer only: the counted proof cites this statement by \textbackslash{}ref as a standing input; the deduction the counted proof has to make is carried out inside the counted span, and no part of that deduction is delegated to this pointer.\par\medskip
\ifdefined\hypertarget\hypertarget{sourcepointer-7}{}\fi
\noindent\textit{Source reference.} Original reference 02KB: focused-seeds/geometry/S9-smg-strong-splitting-affine-scheme/ref-tags.txt, source lines 10--10. Navigation pointer only: the counted proof cites this statement by \textbackslash{}ref as a standing input; the deduction the counted proof has to make is carried out inside the counted span, and no part of that deduction is delegated to this pointer.\par\medskip
\ifdefined\hypertarget\hypertarget{sourcepointer-8}{}\fi
\noindent\textit{Source reference.} Original reference 01ZY: focused-seeds/geometry/S9-smg-strong-splitting-affine-scheme/ref-tags.txt, source lines 11--11. Navigation pointer only: the counted proof cites this statement by \textbackslash{}ref as a standing input; the deduction the counted proof has to make is carried out inside the counted span, and no part of that deduction is delegated to this pointer.\par\medskip
\ifdefined\hypertarget\hypertarget{sourcepointer-9}{}\fi
\noindent\textit{Source reference.} Original reference 01K4: focused-seeds/geometry/S9-smg-strong-splitting-affine-scheme/ref-tags.txt, source lines 12--12. Navigation pointer only: the counted proof cites this statement by \textbackslash{}ref as a standing input; the deduction the counted proof has to make is carried out inside the counted span, and no part of that deduction is delegated to this pointer.\par\medskip
\ifdefined\hypertarget\hypertarget{sourcepointer-10}{}\fi
\noindent\textit{Source reference.} Original reference 04RL: sources/stacks-spaces-more-groupoids/pinned-original-source.tex, source lines 2463--2463. Navigation pointer only: the counted proof cites this statement by \textbackslash{}ref as a standing input; the deduction the counted proof has to make is carried out inside the counted span, and no part of that deduction is delegated to this pointer.\par\medskip
\begin{lemma}
\label{lemma-strong-splitting-affine-scheme}
In Situation \ref{situation-strong-splitting} assume in addition that
$s, t$ are flat and locally of finite presentation and that $U$ is affine.
Then there exists an affine scheme $U'$, an \'etale morphism
$U' \to U$, and a point $u' \in U'$ lying over $u$ with
$\kappa(u) = \kappa(u')$ such that the restriction $R' = R|_{U'}$ of
$R$ to $U'$ is strongly split over $u'$.
\end{lemma}

\begin{proof}
Let $U' \to U$ and $u' \in U'$ be the separated \'etale morphism of schemes
we found in Lemma \ref{lemma-strong-splitting-scheme}.
Let $P \subset R'$ be the strong splitting of $R'$ over $u'$. By
More on Groupoids, Lemma \ref{more-groupoids-lemma-restrict-preserves-type}
the morphisms $s', t' : R' \to U'$ are flat and locally of finite presentation.
They are finite by assumption. Hence $s', t'$ are finite locally
free, see
Morphisms, Lemma \ref{morphisms-lemma-finite-flat}.
In particular $t(s^{-1}(u'))$ is a finite set of points
$\{u'_1, u'_2, \ldots, u'_n\}$ of $U'$. Choose a quasi-compact open
$W \subset U'$ containing each $u'_i$. As $U$ is affine the morphism
$W \to U$ is quasi-compact (see
Schemes, Lemma \ref{schemes-lemma-quasi-compact-affine}).
The morphism $W \to U$ is also locally quasi-finite (see
Morphisms, Lemma \ref{morphisms-lemma-etale-locally-quasi-finite})
and separated. Hence by
More on Morphisms,
Lemma \ref{more-morphisms-lemma-quasi-finite-separated-quasi-affine}
(a version of Zariski's Main Theorem)
we conclude that $W$ is quasi-affine. By
Properties, Lemma \ref{properties-lemma-ample-finite-set-in-affine}
we see that $\{u'_1, \ldots, u'_n\}$ are contained in an affine
open of $U'$. Thus we may apply
Groupoids, Lemma \ref{groupoids-lemma-find-invariant-affine}
to conclude that there exists an affine $P$-invariant open
$U'' \subset U'$ which contains $u'$.

\medskip\noindent
To finish the proof denote $R'' = R|_{U''}$ the restriction of $R$
to $U''$. This is the same as the restriction of $R'$ to $U''$.
As $P \subset R'$ is an open and closed subscheme, so is
$P|_{U''} \subset R''$. By construction the open subscheme $U'' \subset U'$
is $P$-invariant which means that
$P|_{U''} = (s'|_P)^{-1}(U'') = (t'|_P)^{-1}(U'')$
(see discussion in
Groupoids, Section \ref{groupoids-section-invariant})
so the restrictions of $s''$ and $t''$ to $P|_{U''}$ are still finite.
The sub groupoid scheme $P|_{U''}$ is still a strong splitting of
$R''$ over $u''$; above we verified (a), (b) and (c) holds as
$\{r' \in R': t'(r') = u', s'(r') = u'\} =
\{r'' \in R'': t''(r'') = u', s''(r'') = u'\}$ trivially.
The lemma is proved.
\end{proof}

\end{document}
